\section{下一阶段的研究议程：从更强模型到更好语言系统}

\subsection{Benchmark 之后，我们还需要什么}

最后一讲隐含的研究判断是：NLP 不应停在``再追一个更高分数''。当 foundation model 已经把许多通用 benchmark 做到很高以后，研究重点必须继续上移。

\begin{itemize}
    \item \textbf{更细粒度的语言现象评测}：不是只看平均分，而是看模型在指代、蕴含、歧义、语用和风格变化上的表现。
    \item \textbf{更真实的人机任务}：从单轮问答转向长上下文阅读、写作、检索、工具使用与协作任务。
    \item \textbf{更可信的失败分析}：不只展示成功案例，而是系统记录模型在何种语言现象、何种社会语境下容易失败。
\end{itemize}

\begin{importantbox}{NLP 已经从 task solving 进入 system building}
早年的 NLP 论文经常围绕某个清晰任务展开，例如 parsing、NER、sentiment。今天更关键的问题变成了：如何把语言模型做成一个能在真实世界稳定工作的系统。这要求评测、交互、鲁棒性和治理一起前进。
\end{importantbox}

\subsection{grounding 会把很多旧问题重新带回来}

如果语言只在文本内部流转，统计模式就足以带来惊人的效果；但一旦语言要稳定地指向对象、行动和真实环境，问题就会变难。

\begin{enumerate}
    \item \textbf{指称稳定性}：同一句话在不同环境中到底指向什么
    \item \textbf{世界更新}：外部事实变化后，模型如何及时纠正旧知识
    \item \textbf{行动后果}：当模型驱动搜索、代码或 agent 行动时，语言错误会转化为真实代价
\end{enumerate}

\begin{knowledgebox}{为什么 grounded NLP 让语言学重新变得锋利}
一旦语言需要连接视觉、动作、工具和社会情境，语义、语用、篇章与指代这些传统语言学问题就会重新变成系统瓶颈。也正因此，Manning 才强调语言学不是旧时代遗产，而是未来系统设计的分析工具。
\end{knowledgebox}

\subsection{多语言与社会语言学将重新进入主线}

真正全球可用的 NLP 系统，不能只在标准英语上工作。未来研究至少要同时面对三类差异：

\begin{itemize}
    \item 语言之间的资源差异：高资源语种和低资源语种的训练条件截然不同
    \item 同一语言内部的变体差异：方言、口语、社媒语言和专业语体都可能偏离标准书面语
    \item 社会偏见与表示不均：谁的数据更多、谁的话语方式更被模型偏好，都会影响最终系统行为
\end{itemize}

\begin{warningbox}{只在主流英语基准上进步，会制造能力错觉}
一个系统在英文 benchmark 上持续上升，并不意味着它对语言本身理解得更好了。一旦切到低资源语言、口语、方言或强语境文本，很多能力会迅速崩塌。这不是边缘问题，而是 NLP 是否具有普适性的核心检验。
\end{warningbox}

\subsection{本章小结}

后 LLM 时代的 NLP 研究，不应只问模型还能多大、多强，而应更具体地问：它是否真正更可解释、更可泛化、更可治理，也是否更能处理真实世界里复杂、多语言、强语境的语言任务。

